\section{Prices for packets and marked connected components}
\label{sec:mixed-activity-prices}
The vertex bound used in Section~\ref{sec:cluster-expansion} combines
two populations: mandatory defect packets and optional local activities.
We prove the counting and reserve estimates for this combination, then
show how a single insertion is added to its actual alphabet. The analytic
bound on the weight of one packet remains a separate input.

\subsection{Counting physical defect packets}
Fix a block ratio $L\ge3$ and nonnegative integers $\ell,B$. A seed is a fine plaquette, a coarse
oriented bond, or a coarse site. These are distinct types even when
they have the same location. Each seed has a coarse anchor and a
finite set of points within max-distance one of its anchor. For a
plaquette these points are the coarse parents of its touched fine
sites; for a bond they are its two endpoints; for a site it is the
site itself. For the plaquette, the parent-distance estimate is
$L\,\operatorname{dist}(\operatorname{parent}(x),\operatorname{parent}(y))\le L+2$.
Since the distance is integral and $L\ge3$, it is at most one. The halo of radius $\ell B$ around these points is its
support, contained in the ball of radius $\ell B+1$ about the anchor.
A packet is a nonempty finite set of seeds whose halo intersection
graph is connected. Its support is the union of these halos, and
its size $n$ is the number of typed seeds.

The number of seeds with anchors in a set $Y$ is at most $T_L|Y|$,
where
\[
        T_L=L^4\#\{\text{plaquette axes}\}+4+1.
\]
The fine site is determined by its coarse parent and one of $L^4$
residues, which proves the plaquette count; the two other counts are
immediate. Consequently at most
\[
 H=(2(\ell B+1)+1)^4T_L
\]
seed halos can contain a given site. If two halos overlap, their
anchors have distance at most $2(\ell B+1)$, so each seed has at most
\[
                   D=(4(\ell B+1)+1)^4T_L\ge1
\]
possible neighboring seeds. These bounds remain valid when balls wrap
around the periodic lattice, since wrapping only decreases the number
of distinct sites.

There are at most $HD^{2n}$ packets of size $n$ containing a fixed
site in their support. To prove this, choose a seed whose halo contains
that site; there are at most $H$ choices. A connected $n$-seed packet
has a spanning tree. Traversing every tree edge out and back gives a
walk of length $2(n-1)$, starting at the chosen seed, whose visited
set is exactly the packet. There are at most $D^{2(n-1)}$ such walks.
Different visited sets cannot come from the same walk, so this counts
packets without an injectivity assumption on their placements.
Enlarging the bound to $HD^{2n}$ also covers $n=1$ uniformly. There
is no packet with $n=0$.

A single halo has $B$-cover number at most
$s_\ell=(2\ell+3)^4$. For $B\ge1$ it lies in the ball of radius
$(\ell+1)B$, tiled by $B$-balls with centers on the $B$-spaced grid;
there are at most $(2\ell+3)^4$ centers, forming a connected grid.
For $B=0$ use at most $3^4\le s_\ell$ singleton centers. Covering
subadditivity along a spanning tree therefore bounds a packet's
cover number by $ns_\ell$.

Suppose its nonnegative weight obeys the analytic estimate
\begin{equation}\label{eq:mixed-packet-analytic}
                  p_A\le C\exp((s-g)n(A)),\qquad C\ge0,
\end{equation}
where $s$ is the expense per seed and $g$ the gain per seed. For a
support exponent $h\ge0$, put
\[
                   q_p=D^2\exp(s-g+hs_\ell).
\]
If $q_p\le r_p<1$, multiplication by the support weight, the packet
count, and a geometric sum give
\begin{equation}\label{eq:mixed-packet-hit}
 \sum_{\substack{A:\,x\in X_A\\n(A)\ge n_0}}
                  p_Ae^{h c(X_A)}
          \le HC\frac{r_p^{\max(1,n_0)}}{1-r_p}.
\end{equation}
Indeed the summand at fixed size $n$ is at most $HCq_p^n$.
The lower limit is $\max(1,n_0)$ because packets are nonempty.
This explains the single-seed extraction in the packet price.
If a packet has several analytic variants, the hypothesis
\eqref{eq:mixed-packet-analytic} must bound their \emph{sum} at that
packet. It cannot be applied separately to every variant for free.
A variant with smaller support has no larger support price and is
handled by the same fiberwise summation.

\subsection{Reserving the cost of a whole component}
Let $f_i\ge0$ be weights on a combined alphabet, with supports $X_i$,
extra nonnegative costs $d_i$, and nonnegative marks $e_i$. Define
\[
 \widetilde f_i=f_i\exp(2\varepsilon d_i+e_i),\qquad
 D(A)=\sum_{i\in A}(c(X_i)+d_i),\qquad E(A)=\sum_{i\in A}e_i,
\]
where $\varepsilon>0$. Assume the vertex condition
\eqref{eq:cluster-vertex-input} holds for $\widetilde f$ with support
exponent $a+2\varepsilon$, site price $\gamma$, tree reserve $\delta$
and tail ratio $q$. We will use the same tree bound for distinct
connected sets.

For each connected finite set $A$ containing a root label $i$, choose
a spanning tree rooted at $i$. If $|A|-1=m$, the $m!$ enumerations
of its other labels give $m!$ different numbered trees. Their weights
in \eqref{eq:cluster-forest-bound} each contain $1/m!$. Thus the
sum over distinct sets and their chosen trees is bounded by the full
numbered tree sum. Choose a root whose support contains $x$, or
bound by the sum over all such roots. The rooted bound established
in Section~\ref{sec:cluster-expansion}, followed by the geometric
sum in $m$, proves
\begin{equation}\label{eq:mixed-distinct-tail}
 \sum_{\substack{A\ {\rm connected},\ x\in\cup_{i\in A}X_i\\|A|-1\ge m}}
                   \prod_{i\in A}\widetilde f_i
                                 e^{(a+2\varepsilon)c(X_i)}
                        \le\frac{\gamma q^m}{1-q}.
\end{equation}
No logarithmic coefficient is present in this distinct-set estimate.

Let $k,m\ge0$ be integers and restrict the sets further by
$D(A)\ge T$ and $E(A)\ge U$, where $T,U$ are real. The elementary
polynomial estimate is
\[
 (1+D)^k\le P_{\varepsilon,k}e^{\varepsilon D},\qquad
              P_{\varepsilon,k}=\varepsilon^{-k}k!e^\varepsilon
                         \quad(D\ge0).
\]
Since $D\ge T$ and $E\ge U$,
\[
 e^{\varepsilon D}
       \le e^{-\varepsilon T-U}e^{2\varepsilon D+E}.
\]
For connected $A$ we also have $c(\cup_iX_i)\le\sum_ic(X_i)$.
Multiplying these inequalities, and then using
\eqref{eq:mixed-distinct-tail}, proves
\begin{equation}\label{eq:mixed-reserved-tail}
 \sum_A\left(\prod_{i\in A}f_i\right)(1+D(A))^k
                                      e^{a c(\cup_{i\in A}X_i)}
 \le P_{\varepsilon,k}e^{-\varepsilon T-U}
                                      \frac{\gamma q^m}{1-q}.
\end{equation}
The sum on the left has all the preceding connectedness, site,
cardinality, cost and mark restrictions. The two reserves have
separate uses: one pays the polynomial, and the other leaves the
explicit tail $e^{-\varepsilon T}$.

\subsection{Combining the two vertex populations}
For an optional label take its original weight $w_i$, extra cost
$d_i\ge0$ and no mark. Assume its site estimate with the full head
exponent $h=a+2\varepsilon+\delta$ is
\[
        \sup_x\sum_{i:x\in X_i}
             w_i e^{2\varepsilon d_i}e^{h c(X_i)}\le\Gamma_o.
\]
For a packet of size $n$, set extra cost $bn$ and mark $\mu n$,
where $b,\mu\ge0$. Its modified weight is
$p_Ae^{(2\varepsilon b+\mu)n}$. Thus
\eqref{eq:mixed-packet-hit} applies with seed expense replaced by
$s+2\varepsilon b+\mu$, leaving the analytic packet bound
\eqref{eq:mixed-packet-analytic} unchanged. If the resulting packet
site price is $\Gamma_p$, the combined alphabet has site price at
most $\Gamma_p+\Gamma_o$. Choose $\gamma>0$ with
\[
               \Gamma_p+\Gamma_o\le\gamma,
                       \qquad 4h_B\gamma/\delta\le q<1.
\]
This is exactly the vertex condition needed for
\eqref{eq:mixed-reserved-tail}.

If the allowed component contains at least $r$ seeds in total,
then $E(A)\ge\mu r$. Hence the right-hand side becomes
\begin{equation}\label{eq:mixed-ceiling}
       P_{\varepsilon,k}e^{-\varepsilon T-\mu r}
                                \frac{\gamma q^m}{1-q}.
\end{equation}
This bound controls simultaneously excess label count $m$, polynomial
order $k$, minimum seed count $r$, and minimum total cost $T$.
In particular, $r=1$ may be used only for components known to contain
a mandatory packet. The support of an output may be smaller than
the component support. Summing all output majorants over each
component fiber first, a bound by $C_0$ times its component price
gives $C_0$ times \eqref{eq:mixed-ceiling}, by support monotonicity
and nonnegative summation. This is the required rule for summing
analytic variants without losing their multiplicities.

\subsection{Adding a normalized elementary insertion}
The insertion changes one optional label already present in the
alphabet. Its fine support has cover number $c_f$, support exponent
$\nu_f$, and original scalar majorant $w\ge0$. Let $g>0$ denote
the normalization gain and set
\[
             j=g e^{-(\nu_f-\varepsilon)c_f},\qquad
             D=\frac{w e^{\nu_f c_f}e^{-\varepsilon c_f}}g.
\]
Then $Dj=w$ exactly. Thus $D$ is the insertion debt and $j$ the
normalized vertex weight, not two independent prices.

Let $X_*$ be this label's actual coarse footprint, and suppose its
projection geometry gives
\begin{equation}\label{eq:mixed-marked-footprint}
             h c(X_*)\le hP+\varepsilon c_f,
                       \qquad \nu_f\ge4\varepsilon,
\end{equation}
where $h=a+2\varepsilon+\delta$ is the head exponent and $P$ the
fixed halo constant. Its extra cost is $c_f$. Therefore
\[
 j e^{2\varepsilon c_f}e^{h c(X_*)}
    \le g\exp\{hP+(4\varepsilon-\nu_f)c_f\}
    \le g e^{hP}=:R_*.
\]
At a fixed site there is just this one modified label. Replacing its
old nonnegative weight by $j$ thus increases the optional site bound
by at most $R_*$. This follows even if $j$ is smaller than the old
weight, by bounding the new array by the old array plus the single
new entry. Set the new optional price and total vertex price to
\[
                   \Gamma_o'=\Gamma_o+R_*,\qquad
                                  \gamma'=\gamma+R_*.
\]
The packet estimates are unchanged. Once
$4h_B\gamma'/\delta\le q'<1$, the preceding mixed-component bound
applies to the modified alphabet with these prices. Retaining only
components containing the marked label decreases each nonnegative
sum. This produces the marked-molecule estimate used for the
order-zero tail in Section~\ref{sec:leaf-expansion}.

For the actual projected footprint the sufficient geometric condition
is $h(10P/L)\le\varepsilon$, together with the comparison of fine
and coarse cover radii. Indeed the projected cover estimate is
$c(X_*)\le P+(10P/L)c_f$, which gives
\eqref{eq:mixed-marked-footprint}. The fixed block ratio is chosen
large enough for this condition before the coupling is increased.
If $g$ and $\gamma$ have negative upper exponent $q_*<0$ in the
fine coupling, then $R_*=g e^{hP}$ and $\gamma'$ have the same upper
exponent, with $L$ and the halo constants fixed. The factor $h_B$
has upper exponent zero. The degree price therefore has upper
exponent $q_*$, and is eventually bounded by the strictly larger
power $H^{q_*/2}$, which tends to zero. This proves affordability
of the augmented tree price under precisely these rate and geometric
inputs.

The combinatorial part of the mixed and marked estimates is now
proved. The remaining analytic input is the packet weight bound
\eqref{eq:mixed-packet-analytic}, including all its variants, together
with the optional activity site bound and the quoted projection and
rate estimates on the actual selected block step.
